% Generated draft: requires mathematical and source review.
\documentclass[11pt]{article}
\usepackage[margin=25mm]{geometry}
\usepackage{fontspec}
\setmainfont{Latin Modern Roman}
\usepackage{amsmath,amssymb,mathtools,unicode-math}
\setmathfont{Latin Modern Math}
\usepackage{enumitem,needspace,xurl,hyperref}
\setlength{\parindent}{0pt}
\setlength{\parskip}{5pt}
\setlength{\emergencystretch}{3em}
\newcommand{\N}{\mathbb{N}}
\newcommand{\Z}{\mathbb{Z}}
\newcommand{\Q}{\mathbb{Q}}
\newcommand{\R}{\mathbb{R}}
\newcommand{\C}{\mathbb{C}}
\newcommand{\F}{\mathbb{F}}
\newcommand{\A}{\mathbb{A}}
\newcommand{\PP}{\mathbb{P}}
\newcommand{\E}{\mathbb{E}}
\newcommand{\eps}{\varepsilon}
\newcommand{\dd}{\,\mathrm{d}}
\begin{document}
\section*{UnsolvedMath ID: 30000487}
\textbf{Source problem code:} catalogue entry
\subsection{Definitions and mathematical statement}
Retain \(\Omega G\) for the Iwasawa algebra in the source. In the usual mod-\(p\) convention, for a profinite group \(G\) and a prime \(p\),
\[
 \Omega G=\F_p[[G]]=\varprojlim_{U\trianglelefteq_{\mathrm{open}}G}\F_p[G/U].
\]
This convention is distinct from \(\Z_p[[G]]\); the extracted sentence does not explicitly define its coefficient ring or hypotheses on \(G\). Let \(J(\Omega G)\) denote its Jacobson radical, the intersection of the maximal left ideals, which is a two-sided ideal. In the customary pro-\(p\) setting it is the augmentation ideal, but the required group hypotheses are not printed here.

The question for the intended class of groups is whether
\[
 \operatorname{Spec}_{\mathrm{two\text{-}sided}}(\Omega G)
 =\{(0),J(\Omega G)\}.
\]
A prime here is a proper two-sided ideal \(P\) such that, for two-sided ideals \(I_1,I_2\), \(I_1I_2\subseteq P\) implies \(I_1\subseteq P\) or \(I_2\subseteq P\). Thus the equality asks both that these two listed ideals be prime and that there be no others. It is not just a statement about primitive ideals, height-one ideals or completely prime ideals. In particular, primeness of \((0)\) is not silently assumed from the name Iwasawa algebra.

The literature assessment calls this the general almost-simple question, but does not define almost-simple or state compact \(p\)-adic analytic, pro-\(p\), torsion or related conditions. Those conditions are indispensable to the intended scope; the question is not expanded into a universal assertion for arbitrary profinite or abelian groups. The intended coefficient field and whether prime ideals are algebraic or restricted to a specified topological class also need confirmation.

The literature triage checked 2026-08-21 reports that for the relevant congruence subgroups of \(SL_2(\Z_p)\), the zero ideal and unique maximal ideal are the only prime ideals, while the general almost-simple case remains open. The congruence level, allowed primes and hypotheses of that theorem are not given, so the draft does not assert this for every congruence subgroup or every prime. The partially-solved label and editorial verification note remain dated source metadata, not independently checked current conclusions.
\subsection{Short English statement}
For the intended almost-simple group setting, are zero and the Jacobson radical exactly the two-sided prime ideals of the mod-\(p\) completed group algebra? The source does not state the essential hypotheses on the group. Its reported congruence-subgroup case is retained with unspecified level and prime restrictions, rather than generalized beyond the supplied evidence.
\subsection{Sources}
\begin{itemize}
\item Local UnsolvedMath registry, permanent ID 30000487, source code catalogue entry, ``Prime Ideals in Iwasawa Algebras,'' including original and cleaned two-ideal question and the almost-simple/congruence-subgroup assessment.
\item \emph{Pro-\(p\) Extensions of Global Fields and pro-\(p\) Groups (2007)}, Oberwolfach Reports, pp.\ 1463--1536, DOI 10.4171/owr/2006/25: \url{https://doi.org/10.4171/owr/2006/25}. The citation uses 2007, while the DOI and proposed-year metadata use 2006.
\item Supplied reference ``Ardakov--Wei--Zhang'': \url{https://arxiv.org/abs/0710.0624}.
\item Supplied reference ``current survey stating the remaining open case'': \url{https://people.maths.ox.ac.uk/ardakov/IwasawaSurvey.pdf}. The URL is cited as supplied; no PDF has been opened.
\item Literature triage checked 2026-08-21; record updated 2026-09-16. No external source has been visited or independently checked here.
\end{itemize}
\end{document}
